\documentclass[10pt,a4paper]{amsart}
\usepackage{amsmath,amssymb,amsthm,graphicx}
\usepackage{verbatim}
\newenvironment{reference}{\comment}{\endcomment}
%\newenvironment{reference}{}{}
\newenvironment{slogan}{\comment}{\endcomment}
\newenvironment{history}{\comment}{\endcomment}
% For commutative diagrams we use Xy-pic
\usepackage[all]{xy}

% We use 2cell for 2-commutative diagrams.
\xyoption{2cell}
% Theorem environments.
%
\theoremstyle{plain}
\newtheorem{theorem}[subsection]{Theorem}
\newtheorem{proposition}[subsection]{Proposition}
\newtheorem{lemma}[subsection]{Lemma}

\theoremstyle{definition}
\newtheorem{definition}[subsection]{Definition}
\newtheorem{example}[subsection]{Example}
\newtheorem{exercise}[subsection]{Exercise}
\newtheorem{situation}[subsection]{Situation}

\theoremstyle{remark}
\newtheorem{remark}[subsection]{Remark}
\newtheorem{remarks}[subsection]{Remarks}

\numberwithin{equation}{subsection}

% Macros
%
\def\lim{\mathop{\mathrm{lim}}\nolimits}
\def\colim{\mathop{\mathrm{colim}}\nolimits}
\def\Spec{\mathop{\mathrm{Spec}}}
\def\Hom{\mathop{\mathrm{Hom}}\nolimits}
\def\Ext{\mathop{\mathrm{Ext}}\nolimits}
\def\SheafHom{\mathop{\mathcal{H}\!\mathit{om}}\nolimits}
\def\SheafExt{\mathop{\mathcal{E}\!\mathit{xt}}\nolimits}
\def\Sch{\mathit{Sch}}
\def\Mor{\mathop{\mathrm{Mor}}\nolimits}
\def\Ob{\mathop{\mathrm{Ob}}\nolimits}
\def\Sh{\mathop{\mathit{Sh}}\nolimits}
\def\NL{\mathop{N\!L}\nolimits}
\def\CH{\mathop{\mathrm{CH}}\nolimits}
\def\proetale{{pro\text{-}\acute{e}tale}}
\def\etale{{\acute{e}tale}}
\def\QCoh{\mathit{QCoh}}
\def\Ker{\mathop{\mathrm{Ker}}}
\def\Im{\mathop{\mathrm{Im}}}
\def\Coker{\mathop{\mathrm{Coker}}}
\def\Coim{\mathop{\mathrm{Coim}}}

% Boxtimes
%
\DeclareMathSymbol{\boxtimes}{\mathbin}{AMSa}{"02}

%
% Macros for moduli stacks/spaces
%
\def\QCohstack{\mathcal{QC}\!\mathit{oh}}
\def\Cohstack{\mathcal{C}\!\mathit{oh}}
\def\Spacesstack{\mathcal{S}\!\mathit{paces}}
\def\Quotfunctor{\mathrm{Quot}}
\def\Hilbfunctor{\mathrm{Hilb}}
\def\Curvesstack{\mathcal{C}\!\mathit{urves}}
\def\Polarizedstack{\mathcal{P}\!\mathit{olarized}}
\def\Complexesstack{\mathcal{C}\!\mathit{omplexes}}
% \Pic is the operator that assigns to X its picard group, usage \Pic(X)
% \Picardstack_{X/B} denotes the Picard stack of X over B
% \Picardfunctor_{X/B} denotes the Picard functor of X over B
\def\Pic{\mathop{\mathrm{Pic}}\nolimits}
\def\Picardstack{\mathcal{P}\!\mathit{ic}}
\def\Picardfunctor{\mathrm{Pic}}
\def\Deformationcategory{\mathcal{D}\!\mathit{ef}}
\usepackage{cleveref}
\crefname{originalreference}{original reference}{original references}
\Crefname{originalreference}{Original reference}{Original references}
\makeatletter
\newlabel{spaces-morphisms-lemma-locally-quasi-finite}{{spaces-morphisms-lemma-locally-quasi-finite}{source}{}{sourcepointer-1}{}}
\newlabel{spaces-morphisms-lemma-locally-quasi-finite@cref}{{[originalreference][0][]spaces-morphisms-lemma-locally-quasi-finite}{[1][0][]source}{}{}{}}
\makeatother
\crefname{originalreference}{original reference}{original references}
\Crefname{originalreference}{Original reference}{Original references}
\makeatletter
\newlabel{stacks-properties-lemma-check-property-covering}{{stacks-properties-lemma-check-property-covering}{source}{}{sourcepointer-2}{}}
\newlabel{stacks-properties-lemma-check-property-covering@cref}{{[originalreference][0][]stacks-properties-lemma-check-property-covering}{[1][0][]source}{}{}{}}
\makeatother
\crefname{originalreference}{original reference}{original references}
\Crefname{originalreference}{Original reference}{Original references}
\makeatletter
\newlabel{stacks-properties-section-properties-morphisms}{{stacks-properties-section-properties-morphisms}{source}{}{sourcepointer-3}{}}
\newlabel{stacks-properties-section-properties-morphisms@cref}{{[originalreference][0][]stacks-properties-section-properties-morphisms}{[1][0][]source}{}{}{}}
\makeatother
\title{Local quasi-finiteness of a representable morphism is equivalent to discreteness of all field fibres}
\author{The Stacks Project Authors}
\date{stacks-project:stacks-morphisms.tex:local-snapshot-20261012; GNU FDL 1.2 or later}
\begin{document}
\maketitle
\noindent\textit{Editorial note.} All delivered mathematical text is the author's own, copied from the pinned Stacks Project chapter source sources/stacks-algebraic-stacks-morphisms/pinned-original-source.tex; no statement or proof step was rewritten, reordered or removed.\par\medskip
\noindent\textit{Editorial note.} Preamble provenance: the Stacks Project keeps one single shared preamble for every chapter; this card takes it directly from the staged sources/stacks-algebraic-stacks-morphisms/preamble.tex as three verbatim spans (lines 8-12, 14-18 and 154-217), so no derived file was written and nothing was added or redefined. The wrapper class and the base packages are supplied by the assembler.\par\medskip
\noindent\textit{Editorial note.} Bibliography provenance: neither the counted statement nor its proof contains a citation (this lane rescans the whole span before writing the card), so the card carries no bibliography and no entry was reconstructed.\par\medskip
\noindent\textit{Editorial note.} Selection provenance: this card was opened under the file-based control contract (CONTROL.md), with the mechanical pre-card gate run on its statement and proof interval before the card was written.\par\medskip
\noindent\textit{Allowed standard prerequisites (editorial).}\begin{itemize}
\item The statement that a locally quasi-finite morphism of algebraic spaces has discrete fibres over fields. Invoked inside the counted proof (source lines 4929-4932) as a cross-chapter pointer. Gives discreteness of the fibre over a field from local quasi-finiteness.
\item The criterion for checking a property of morphisms of algebraic stacks along a smooth surjective cover. Invoked inside the counted proof (source lines 4936-4940) as a cross-chapter pointer. Reduces local quasi-finiteness of the morphism of stacks to that of a base change covering it.
\item The characterization of local quasi-finiteness of a morphism of algebraic spaces by local finite type plus discreteness of field fibres. Invoked at the end of the counted proof (source lines 4945-4948) as a cross-chapter pointer. Concludes local quasi-finiteness of the base-changed morphism.
\end{itemize}\medskip
\ifdefined\hypertarget\hypertarget{sourcepointer-1}{}\fi
\noindent\textit{Source reference.} Original reference spaces-morphisms-lemma-locally-quasi-finite: sources/stacks-algebraic-stacks-morphisms/pinned-original-source.tex, source lines 4931--4931. Navigation pointer only: the counted text refers to the Stacks item spaces-morphisms-lemma-locally-quasi-finite, which belongs to another chapter of the project and is not part of this card. The pointer carries the target's own label as its number because this chapter's staged snapshot does not contain the other chapter's tag table; its content is not reproduced here.\par\medskip
\ifdefined\hypertarget\hypertarget{sourcepointer-2}{}\fi
\noindent\textit{Source reference.} Original reference stacks-properties-lemma-check-property-covering: sources/stacks-algebraic-stacks-morphisms/pinned-original-source.tex, source lines 479--479. Navigation pointer only: the counted text refers to the Stacks item stacks-properties-lemma-check-property-covering, which belongs to another chapter of the project and is not part of this card. The pointer carries the target's own label as its number because this chapter's staged snapshot does not contain the other chapter's tag table; its content is not reproduced here.\par\medskip
\ifdefined\hypertarget\hypertarget{sourcepointer-3}{}\fi
\noindent\textit{Source reference.} Original reference stacks-properties-section-properties-morphisms: sources/stacks-algebraic-stacks-morphisms/pinned-original-source.tex, source lines 248--248. Navigation pointer only: the counted text refers to the Stacks item stacks-properties-section-properties-morphisms, which belongs to another chapter of the project and is not part of this card. The pointer carries the target's own label as its number because this chapter's staged snapshot does not contain the other chapter's tag table; its content is not reproduced here.\par\medskip
\begin{lemma}
\label{lemma-representable-by-spaces-quasi-finite}
Let $f : \mathcal{X} \to \mathcal{Y}$ be a morphism of algebraic stacks.
Assume $f$ is representable by algebraic spaces.
The following are equivalent
\begin{enumerate}
\item $f$ is locally quasi-finite (as in Properties of Stacks,
Section \ref{stacks-properties-section-properties-morphisms}), and
\item $f$ is locally of finite type and for every morphism
$\Spec(k) \to \mathcal{Y}$ where $k$ is a field the
space $|\Spec(k) \times_\mathcal{Y} \mathcal{X}|$ is discrete.
\end{enumerate}
\end{lemma}

\begin{proof}
Assume (1). In this case the morphism of algebraic spaces
$\mathcal{X}_k \to \Spec(k)$ is locally quasi-finite as a base change
of $f$. Hence $|\mathcal{X}_k|$ is discrete by
Morphisms of Spaces, Lemma \ref{spaces-morphisms-lemma-locally-quasi-finite}.
Conversely, assume (2). Pick a surjective smooth morphism
$V \to \mathcal{Y}$ where $V$ is a scheme. It suffices to show that the
morphism of algebraic spaces $V \times_\mathcal{Y} \mathcal{X} \to V$
is locally quasi-finite, see
Properties of Stacks, Lemma
\ref{stacks-properties-lemma-check-property-covering}.
The morphism $V \times_\mathcal{Y} \mathcal{X} \to V$ is locally of finite
type by assumption. For any morphism $\Spec(k) \to V$ where $k$ is a
field
$$
\Spec(k) \times_V (V \times_\mathcal{Y} \mathcal{X}) =
\Spec(k) \times_\mathcal{Y} \mathcal{X}
$$
has a discrete space of points by assumption. Hence we conclude that
$V \times_\mathcal{Y} \mathcal{X} \to V$ is locally quasi-finite by
Morphisms of Spaces, Lemma \ref{spaces-morphisms-lemma-locally-quasi-finite}.
\end{proof}

\end{document}
